% GENERATED FILE. Edit canonical lessons, then run npm run generate:books.
% canonical-source-hash: fnv1a64:e6401c451d2746f2
% canonical-lessons: GU-C75-ahead, GU-C75-behind, GU-C75-inside, GU-C75-outside, GU-C75-between

\chapter{Ahead, Behind, In and Out}
\label{ch:gu-ahead-behind-in-and-out}
\hlchaptermodality{75}

\begin{chapteropening}
\textbf{By the end of this chapter:} \emph{I can say ahead, behind, inside, outside and between.}

Complete the last lesson of chapter 75: I can say ahead, behind, inside, outside and between.
\end{chapteropening}

\section[\texorpdfstring{\gu{આગળ} (āgaḷ)}{આગળ (āgaḷ)}]{\gu{આગળ} (āgaḷ) — ahead, in front}

\label{lesson:GU-C75-ahead}



\begin{quote}
Before the new one: say the Gujarati for up, then the Gujarati for down.
\end{quote}

\subsection*{You'll want to know: \gu{આગળ}}
\textbf{\gu{આગળ}} (\emph{āgaḷ}) — \textquotedblleft{}ahead, in front\textquotedblright{}.

\begin{cousinweb}
\textbf{\gu{આગળ}} (\emph{āgaḷ}) is a tadbhava. CDIAL lists it under Sanskrit \emph{agra}, \textquotedblleft{}top, front\textquotedblright{}, among the forms built with an \emph{-l-} suffix (CDIAL 68).
\end{cousinweb}

\begin{grammarlens}[title={What you've built}]
One more word for ahead, behind, in and out.
\end{grammarlens}

\subsection*{Your turn}
\begin{itemize}
\raggedright
  \item \emph{Say it:} \emph{āgaḷ}
  \item \emph{Say it:} \emph{āgaḷ}, once more
  \item \emph{Say it:} say \emph{āgaḷ}
  \item \emph{Recall:} say the Gujarati for up, then the Gujarati for down, then say \emph{āgaḷ} again
\end{itemize}

\begin{tcolorbox}[breakable,colback=teal!4,colframe=teal!35!black,title={Before you move on}]
What does \gu{આગળ} mean? (\textquotedblleft{}Ahead, in front\textquotedblright{}.) Say it once more.
\end{tcolorbox}
\section[\texorpdfstring{\gu{પાછળ} (pāchaḷ)}{પાછળ (pāchaḷ)}]{\gu{પાછળ} (pāchaḷ) — behind}

\label{lesson:GU-C75-behind}



\begin{quote}
Before the new one: say the Gujarati for down, then the Gujarati for ahead.
\end{quote}

\subsection*{You'll want to know: \gu{પાછળ}}
\textbf{\gu{પાછળ}} (\emph{pāchaḷ}) — \textquotedblleft{}behind\textquotedblright{}.

\begin{cousinweb}
\textbf{\gu{પાછળ}} (\emph{pāchaḷ}) is built on the stem of \emph{pāchũ}, \textquotedblleft{}back\textquotedblright{}, which CDIAL lists under \emph{*paśca}, \textquotedblleft{}hinder part\textquotedblright{} (CDIAL 7990). CDIAL does not list \textbf{\gu{પાછળ}} (\emph{pāchaḷ}) itself.
\end{cousinweb}

\begin{grammarlens}[title={What you've built}]
One more word for ahead, behind, in and out.
\end{grammarlens}

\subsection*{Your turn}
\begin{itemize}
\raggedright
  \item \emph{Say it:} \emph{pāchaḷ}
  \item \emph{Say it:} \emph{pāchaḷ}, once more
  \item \emph{Say it:} say \emph{pāchaḷ}
  \item \emph{Recall:} say the Gujarati for down, then the Gujarati for ahead, then say \emph{pāchaḷ} again
\end{itemize}

\begin{tcolorbox}[breakable,colback=teal!4,colframe=teal!35!black,title={Before you move on}]
What does \gu{પાછળ} mean? (\textquotedblleft{}Behind\textquotedblright{}.) Say it once more.
\end{tcolorbox}
\section[\texorpdfstring{\gu{અંદર} (andar)}{અંદર (andar)}]{\gu{અંદર} (andar) — inside}

\label{lesson:GU-C75-inside}



\begin{quote}
Before the new one: say the Gujarati for ahead, then the Gujarati for behind.
\end{quote}

\subsection*{You'll want to know: \gu{અંદર}}
\textbf{\gu{અંદર}} (\emph{andar}) — \textquotedblleft{}inside\textquotedblright{}.

\begin{cousinweb}
\textbf{\gu{અંદર}} (\emph{andar}) is a loan, usually traced to Persian \emph{andar}, \textquotedblleft{}inside\textquotedblright{}.
\end{cousinweb}

\begin{grammarlens}[title={What you've built}]
One more word for ahead, behind, in and out.
\end{grammarlens}

\subsection*{Your turn}
\begin{itemize}
\raggedright
  \item \emph{Say it:} \emph{andar}
  \item \emph{Say it:} \emph{andar}, once more
  \item \emph{Say it:} say \emph{andar}
  \item \emph{Recall:} say the Gujarati for ahead, then the Gujarati for behind, then say \emph{andar} again
\end{itemize}

\begin{tcolorbox}[breakable,colback=teal!4,colframe=teal!35!black,title={Before you move on}]
What does \gu{અંદર} mean? (\textquotedblleft{}Inside\textquotedblright{}.) Say it once more.
\end{tcolorbox}
\section[\texorpdfstring{\gu{બહાર} (bahār)}{બહાર (bahār)}]{\gu{બહાર} (bahār) — outside}

\label{lesson:GU-C75-outside}



\begin{quote}
Before the new one: say the Gujarati for behind, then the Gujarati for inside.
\end{quote}

\subsection*{You'll want to know: \gu{બહાર}}
\textbf{\gu{બહાર}} (\emph{bahār}) — \textquotedblleft{}outside\textquotedblright{}.

\begin{cousinweb}
\textbf{\gu{બહાર}} (\emph{bahār}) is a tadbhava, a word inherited and reshaped by sound change. It comes from \emph{*bāhira}, \textquotedblleft{}external\textquotedblright{}, a form CDIAL reconstructs from its descendants (CDIAL 9226). It is built on Sanskrit \emph{bahis}, \textquotedblleft{}outside\textquotedblright{}.
\end{cousinweb}

\begin{grammarlens}[title={What you've built}]
One more word for ahead, behind, in and out.
\end{grammarlens}

\subsection*{Your turn}
\begin{itemize}
\raggedright
  \item \emph{Say it:} \emph{bahār}
  \item \emph{Say it:} \emph{bahār}, once more
  \item \emph{Say it:} say \emph{bahār}
  \item \emph{Recall:} say the Gujarati for behind, then the Gujarati for inside, then say \emph{bahār} again
\end{itemize}

\begin{tcolorbox}[breakable,colback=teal!4,colframe=teal!35!black,title={Before you move on}]
What does \gu{બહાર} mean? (\textquotedblleft{}Outside\textquotedblright{}.) Say it once more.
\end{tcolorbox}
\section[\texorpdfstring{\gu{વચ્ચે} (vacce)}{વચ્ચે (vacce)}]{\gu{વચ્ચે} (vacce) — between, in the middle}

\label{lesson:GU-C75-between}



\begin{quote}
Before the new one: say the Gujarati for inside, then the Gujarati for outside.
\end{quote}

\subsection*{You'll want to know: \gu{વચ્ચે}}
\textbf{\gu{વચ્ચે}} (\emph{vacce}) — \textquotedblleft{}between, in the middle\textquotedblright{}.

\begin{cousinweb}
\textbf{\gu{વચ્ચે}} (\emph{vacce}) is a tadbhava, a word inherited and reshaped by sound change. It comes from \emph{*vīcya}, \textquotedblleft{}middle\textquotedblright{}, a form CDIAL reconstructs from its descendants (CDIAL 12042).
\end{cousinweb}

\begin{grammarlens}[title={What you've built}]
One more word for ahead, behind, in and out.
\end{grammarlens}

\subsection*{Your turn}
\begin{itemize}
\raggedright
  \item \emph{Say it:} \emph{vacce}
  \item \emph{Say it:} \emph{vacce}, once more
  \item \emph{Say it:} say \emph{vacce}
  \item \emph{Recall:} say the Gujarati for inside, then the Gujarati for outside, then say \emph{vacce} again
\end{itemize}

\begin{tcolorbox}[breakable,colback=teal!4,colframe=teal!35!black,title={Before you move on}]
What does \gu{વચ્ચે} mean? (\textquotedblleft{}Between, in the middle\textquotedblright{}.) Say it once more.
\end{tcolorbox}
